\documentclass[11pt,letterpaper]{article}

\usepackage[utf8]{inputenc}
\usepackage[T1]{fontenc}
\usepackage[spanish,es-tabla,es-noshorthands,es-nodecimaldot]{babel}
\usepackage{lmodern}
\usepackage[margin=2.5cm]{geometry}
\usepackage{amsmath,amssymb,amsthm}
\usepackage{booktabs}
\usepackage{array}
\usepackage{multirow}
\usepackage{xcolor}
\usepackage{listings}
\usepackage[spanish,ruled,vlined,linesnumbered]{algorithm2e}
\usepackage{pgfplots}
\pgfplotsset{compat=1.16}
\usepackage{hyperref}

\hypersetup{colorlinks=true,linkcolor=blue!50!black,urlcolor=blue!50!black}

\newtheorem{teorema}{Teorema}
\newtheorem{lema}{Lema}
\newtheorem{proposicion}{Proposición}

\newcommand{\E}[1]{\ensuremath{\langle #1 \rangle}}
\newcommand{\Ext}{\operatorname{Ext}}
\newcommand{\Esf}{\operatorname{Esfuerzo}}
\newcommand{\Mod}{\operatorname{Mod}}
\newcommand{\Rmax}{R_{\max}}

\lstset{
  basicstyle=\ttfamily\small,
  keywordstyle=\color{blue!60!black}\bfseries,
  commentstyle=\color{green!40!black},
  numbers=left,
  numberstyle=\tiny,
  frame=single,
  breaklines=true,
  columns=fullflexible,
  language=C++
}

\title{\textbf{Moderando el extremismo de opiniones en una red social (ModEx)}\\[0.5em]
\large Proyecto de curso -- Análisis y Diseño de Algoritmos II\\
Escuela de Ingeniería de Sistemas y Computación}
% TODO: completar con los nombres y códigos de los integrantes del grupo.
\author{Santiago Arias Rojas (2416285) \and Sebastian Calvo Carvajal (2419918) \and Sebastian Rubio Piedrahita (2459628)\\[0.5em]
Profesor: Jesús Alexander Aranda \quad Monitor: Samuel Galindo}
\date{Noviembre de 2026}

\begin{document}

\maketitle
\tableofcontents
\newpage

%=====================================================================
\section{Introducción}
%=====================================================================

En este informe se estudia el problema ModEx: dada una red social de $n$
agentes, cada uno con una opinión, una receptividad y una influencia, y un
valor máximo $\Rmax$ disponible para moderar opiniones, se debe escoger a qué
agentes moderar (llevar su opinión a $0$) para que el extremismo de la red
resultante sea mínimo sin superar $\Rmax$.

Se diseñaron, analizaron e implementaron tres soluciones: fuerza bruta
(\texttt{modexFB}), un algoritmo voraz (\texttt{modexV}) y programación
dinámica (\texttt{modexPD}). Para cada una se estudia su complejidad y su
corrección, y al final se comparan experimentalmente y se concluye sobre sus
ventajas y desventajas en la práctica.

%=====================================================================
\section{El problema}
%=====================================================================

\subsection{Formalización}

Una red social es una pareja $RS = \langle AG, \Rmax \rangle$ con
$AG = \langle a_0,\dots,a_{n-1}\rangle$ y $\Rmax \in \mathbb{N}$. Cada agente
es una terna $a_i = \langle o_i, r_i, inf_i\rangle$ con
$o_i \in \{-100,\dots,100\}$, $r_i \in [0,1]$ e $inf_i \in [1,3]$.

El extremismo de la red es
\[
  \Ext(RS) = \frac{\sqrt{\sum_{i=0}^{n-1} o_i^2 \cdot inf_i}}{n}.
\]
Esta es la lectura de la fórmula que reproduce los valores de los ejemplos
del enunciado ($47.14$, $33.5$ y $9.43$), y es la que usa la implementación.

Una estrategia es $E = \langle e_0,\dots,e_{n-1}\rangle$ con
$e_i \in \{0,1\}$; $\Mod(RS,E)$ cambia a $0$ la opinión de los agentes con
$e_i = 1$. El esfuerzo es
\[
  \Esf(RS,E) = \sum_{i=0}^{n-1}\left\lceil |o_i - o_i'| \cdot (1-r_i)\cdot inf_i \right\rceil .
\]
$E$ es aplicable si $\Esf(RS,E) \le \Rmax$. ModEx pide una estrategia
aplicable que minimice $\Ext(\Mod(RS,E))$.

\subsection{Reformulación: una mochila 0/1}
\label{sec:mochila}

Para cada agente definimos su \emph{costo} y su \emph{beneficio}:
\[
  c_i = \left\lceil |o_i|\cdot(1-r_i)\cdot inf_i \right\rceil ,
  \qquad
  b_i = o_i^2 \cdot inf_i .
\]
Si $e_i = 0$ entonces $o_i' = o_i$ y el término del esfuerzo es
$\lceil 0 \rceil = 0$; si $e_i = 1$ entonces $o_i' = 0$ y el término es $c_i$.
Por lo tanto
\[
  \Esf(RS,E) = \sum_{i} e_i\, c_i ,
  \qquad
  \Ext(\Mod(RS,E)) = \frac{\sqrt{S - B(E)}}{n},
  \quad\text{con } S = \sum_i b_i,\; B(E) = \sum_i e_i\, b_i .
\]
Como $x \mapsto \sqrt{S-x}/n$ es decreciente, \textbf{minimizar el extremismo
equivale a maximizar $B(E)$ sujeto a $\sum_i e_i c_i \le \Rmax$}: es
exactamente el problema de la mochila 0/1, donde los agentes son los objetos,
$c_i$ su peso, $b_i$ su valor y $\Rmax$ la capacidad. Esta observación guía el
diseño del algoritmo voraz y del de programación dinámica.

Dos cotas útiles: como $|o_i| \le 100$, $1-r_i \le 1$ e $inf_i \le 3$, se
tiene $0 \le c_i \le 300$ y por lo tanto $\sum_i c_i \le 300n$.

Además, la estrategia $E = \langle 0,\dots,0\rangle$ tiene esfuerzo $0 \le
\Rmax$, así que \textbf{siempre existe al menos una estrategia aplicable}.

\subsection{Ejemplos del enunciado}

\begin{table}[h]
\centering
\begin{tabular}{lcccc}
\toprule
Red & Agente & $\langle o, r, inf\rangle$ & $c_i$ & $b_i$ \\
\midrule
\multirow{3}{*}{$RS_1$, $\Rmax = 55$}
 & $a_0$ & $\langle 100, 0.5, 1\rangle$ & 50 & 10000 \\
 & $a_1$ & $\langle 100, 0.1, 1\rangle$ & 90 & 10000 \\
 & $a_2$ & $\langle -10, 0.1, 1\rangle$ & 9  & 100 \\
\midrule
\multirow{3}{*}{$RS_2$, $\Rmax = 35$}
 & $a_0$ & $\langle -30, 0.9, 1\rangle$ & 3  & 900 \\
 & $a_1$ & $\langle 20, 0.1, 2\rangle$  & 36 & 800 \\
 & $a_2$ & $\langle 50, 0.5, 1\rangle$  & 25 & 2500 \\
\bottomrule
\end{tabular}
\caption{Costos y beneficios de los agentes de los ejemplos 2.3.1 y 2.3.2.}
\label{tab:ejemplos}
\end{table}

En $RS_1$ las estrategias aplicables no vacías son $\E{1,0,0}$
($\Ext = 33.50$) y $\E{0,0,1}$ ($\Ext = 47.14$); $\E{1,0,1}$ cuesta $59 > 55$ y
cualquier estrategia que modere a $a_1$ cuesta al menos $90$. El óptimo es
$\E{1,0,0}$ con esfuerzo $50$ y extremismo $33.50$.

En $RS_2$ el óptimo es $\E{1,0,1}$ con esfuerzo $28$ y extremismo
$\sqrt{800}/3 = 9.43$: moderar a $a_1$ cuesta $36 > 35$, así que solo se
pueden moderar subconjuntos de $\{a_0, a_2\}$ y moderar a ambos es lo mejor.

%=====================================================================
\section{Fuerza bruta}
%=====================================================================

El algoritmo genera todas las estrategias posibles y entre ellas escoge la
mejor aplicable. Las estrategias se generan recursivamente fijando $e_i = 0$ y
luego $e_i = 1$ para cada agente; cada estrategia completa se evalúa apenas se
genera, de modo que no es necesario guardar las $2^n$ estrategias en memoria.

\begin{algorithm}[h]
\caption{Fuerza bruta (\texttt{modexFB})}
\SetKwFunction{Explorar}{Explorar}
\SetKwProg{Fn}{Función}{:}{}
\KwIn{Red $RS = \langle AG, \Rmax\rangle$ con $n$ agentes}
\KwOut{$(E, \Esf(RS,E), \Ext(\Mod(RS,E)))$ óptima}
$E \gets \langle 0,\dots,0\rangle$; $\mathit{mejor} \gets (\emptyset, 0, +\infty)$\;
\Explorar{$0$}\;
\Return $\mathit{mejor}$\;
\BlankLine
\Fn{\Explorar{$i$}}{
  \eIf{$i = n$}{
    $\mathit{esf} \gets \Esf(RS, E)$\tcp*{$O(n)$}
    \If{$\mathit{esf} \le \Rmax$}{
      $\mathit{ext} \gets \Ext(\Mod(RS,E))$\tcp*{$O(n)$}
      \If{$\mathit{ext} < \mathit{mejor}.\mathit{ext}$ \textbf{o} ($\mathit{ext} = \mathit{mejor}.\mathit{ext}$ \textbf{y} $\mathit{esf} < \mathit{mejor}.\mathit{esf}$)}{
        $\mathit{mejor} \gets (E, \mathit{esf}, \mathit{ext})$\;
      }
    }
  }{
    $e_i \gets 0$; \Explorar{$i+1$}\;
    $e_i \gets 1$; \Explorar{$i+1$}\;
    $e_i \gets 0$\;
  }
}
\end{algorithm}

\subsection{Complejidad}

El árbol de recursión es un árbol binario completo de altura $n$: tiene $2^n$
hojas (una por estrategia) y $2^n - 1$ nodos internos. Cada nodo interno hace
trabajo constante y cada hoja calcula el esfuerzo y el extremismo en $O(n)$.
Equivalentemente, $T(n)$ cumple $T(i) = 2T(i+1) + O(1)$ para $i < n$ y
$T(n) = O(n)$ en las hojas, de donde
\[
  T(n) = (2^n - 1)\cdot O(1) + 2^n \cdot O(n) = \Theta(n \cdot 2^n).
\]
En espacio, la profundidad de la recursión es $n$ y solo se guardan la
estrategia actual y la mejor, por lo que el espacio es $\Theta(n)$.

Una primera versión del código construía la lista completa de las $2^n$
estrategias antes de evaluarlas, lo que requería $\Theta(n\cdot 2^n)$ de
memoria (unos 190 MB para $n = 20$). Se reemplazó por la exploración
recursiva descrita, que tiene el mismo orden de tiempo y espacio lineal.

\subsection{Corrección}

\textbf{Sí, el algoritmo siempre da la respuesta correcta.} Toda estrategia de
moderación es una secuencia de $n$ bits, y la recursión visita cada una de las
$2^n$ secuencias exactamente una vez (en cada nivel $i$ se toman las dos
opciones de $e_i$). Entre ellas descarta solo las no aplicables y conserva la
de menor extremismo. Como $\E{0,\dots,0}$ siempre es aplicable, el conjunto de
candidatas no es vacío y el resultado es una estrategia aplicable de
extremismo mínimo. En caso de empate en extremismo se prefiere la de menor
esfuerzo, lo que no afecta la optimalidad.

\paragraph{Detalle de implementación: el redondeo.} El esfuerzo se calcula con
números de punto flotante, donde por ejemplo $1 - 0.7 = 0.30000000000000004$.
Para el agente $\langle 10, 0.7, 1\rangle$ el valor exacto es
$\lceil 10\cdot 0.3 \cdot 1\rceil = 3$, pero aplicar \texttt{ceil}
directamente da $4$, con lo que una estrategia aplicable (por ejemplo con
$\Rmax = 3$) se descartaría. La función \texttt{techo} resta una tolerancia de
$10^{-9}$ antes de aplicar \texttt{ceil}, muy inferior a la precisión de los
datos de entrada. Este caso está incluido en las pruebas (caso $E_7$ de la
tabla~\ref{tab:voraz}).

%=====================================================================
\section{Algoritmo voraz}
%=====================================================================

\subsection{Descripción del algoritmo}

Por la reformulación de la sección~\ref{sec:mochila}, cada agente aporta un
beneficio $b_i$ (reducción de $S$) a cambio de un costo $c_i$. La idea voraz
es moderar primero a los agentes que más reducen el extremismo \emph{por
unidad de esfuerzo}, es decir, los de mayor razón $b_i / c_i$:

\begin{enumerate}
  \item Descartar los agentes con $b_i = 0$ (opinión $0$, no aportan) y los
        que cuestan más que $\Rmax$ (nunca pueden moderarse).
  \item Ordenar los agentes restantes por $b_i/c_i$ de mayor a menor. Los de
        costo $0$ (receptividad $1$) tienen razón $+\infty$ y van primero; los
        empates se rompen por mayor beneficio.
  \item Recorrer la lista en ese orden y moderar a cada agente cuyo costo aún
        quepa en el presupuesto restante. El recorrido no se detiene en el
        primer agente que no cabe: sigue buscando agentes más baratos.
  \item Repetir el paso 3 empezando con el agente de mayor beneficio
        $b_{\max}$ ya moderado, y quedarse con la mejor de las dos soluciones.
\end{enumerate}

El paso 4 corrige el caso en que la razón favorece a agentes baratos pero
poco extremos que consumen el presupuesto que necesitaba un agente muy
extremo (caso $E_4$ de la tabla~\ref{tab:voraz}). Cuesta solo un recorrido
adicional y, como se demuestra en la sección~\ref{sec:voraz-correccion}, da
una garantía de calidad que el paso 3 solo no tiene.

\begin{algorithm}[h]
\caption{Algoritmo voraz (\texttt{modexV})}
\SetKwFunction{Llenar}{Llenar}
\SetKwProg{Fn}{Función}{:}{}
\KwIn{Red $RS = \langle AG, \Rmax\rangle$ con $n$ agentes}
\KwOut{$(E, \Esf(RS,E), \Ext(\Mod(RS,E)))$}
$L \gets \{\, i : b_i > 0 \text{ y } c_i \le \Rmax \,\}$\;
ordenar $L$ por $b_i/c_i$ descendente (razón $+\infty$ si $c_i = 0$)\;
$E \gets \langle 0,\dots,0\rangle$; $(E, B) \gets$ \Llenar{$E$, $0$, $0$}\;
$m \gets \arg\max_{i \in L} b_i$\;
\If{$L \neq \emptyset$ \textbf{y} $e_m = 0$}{
  $E' \gets \langle 0,\dots,0\rangle$; $e'_m \gets 1$\;
  $(E', B') \gets$ \Llenar{$E'$, $c_m$, $b_m$}\;
  \lIf{$B' > B$}{$E \gets E'$}
}
\Return $(E, \Esf(RS,E), \Ext(\Mod(RS,E)))$\;
\BlankLine
\Fn{\Llenar{$E$, $\mathit{usado}$, $B$}}{
  \For{$i \in L$ en orden}{
    \If{$e_i = 0$ \textbf{y} $\mathit{usado} + c_i \le \Rmax$}{
      $e_i \gets 1$; $\mathit{usado} \gets \mathit{usado} + c_i$; $B \gets B + b_i$\;
    }
  }
  \Return $(E, B)$\;
}
\end{algorithm}

\subsection{Entendimos el algoritmo}

\paragraph{Ejemplo 2.3.1 ($RS_1$, $\Rmax = 55$).} Las razones son
$a_0: 10000/50 = 200$, $a_1: 10000/90 = 111.1$ y $a_2: 100/9 = 11.1$, pero
$a_1$ se descarta porque $c_1 = 90 > 55$. El orden es $a_0, a_2$. Se modera a
$a_0$ (usado $= 50$); $a_2$ no cabe ($59 > 55$). En el paso 4, el agente de
mayor beneficio es $a_0$, que ya está moderado. La salida es $E = \E{1,0,0}$,
con esfuerzo $50$ y extremismo $\sqrt{10100}/3 = 33.50$. \textbf{Es la
solución óptima.}

\paragraph{Ejemplo 2.3.2 ($RS_2$, $\Rmax = 35$).} Se descarta $a_1$
($c_1 = 36 > 35$). Las razones son $a_0: 900/3 = 300$ y $a_2: 2500/25 = 100$.
Se modera a $a_0$ (usado $= 3$) y luego a $a_2$ (usado $= 28$). En el paso 4,
el de mayor beneficio es $a_2$, ya moderado. La salida es $E = \E{1,0,1}$,
con esfuerzo $28$ y extremismo $\sqrt{800}/3 = 9.43$. \textbf{Es la solución
óptima.}

\paragraph{Otras entradas.} Se diseñaron seis entradas adicionales, cada una
para ejercitar una situación distinta (tabla~\ref{tab:entradas}). La
tabla~\ref{tab:voraz} compara la salida del voraz con el óptimo, calculado con
fuerza bruta y confirmado con programación dinámica.

\begin{table}[h]
\centering
\small
\begin{tabular}{lcp{6.6cm}p{4.4cm}}
\toprule
Caso & $\Rmax$ & Agentes $\langle o, r, inf\rangle$ & Propósito \\
\midrule
$E_1$ & 55  & $\langle 100,0.5,1\rangle,\langle 100,0.1,1\rangle,\langle -10,0.1,1\rangle$ & Ejemplo 2.3.1 \\
$E_2$ & 35  & $\langle -30,0.9,1\rangle,\langle 20,0.1,2\rangle,\langle 50,0.5,1\rangle$ & Ejemplo 2.3.2 \\
$E_3$ & 10  & $\langle 60,0.9,1\rangle,\langle 50,0.9,1\rangle,\langle 50,0.9,1\rangle$ & Contraejemplo del voraz \\
$E_4$ & 50  & $\langle 20,0.95,1\rangle,\langle 100,0.5,1\rangle$ & Necesidad del paso 4 \\
$E_5$ & 0   & $\langle 80,1,2\rangle,\langle -50,0.2,1.5\rangle,\langle 0,0.3,3\rangle,\langle 30,0.6,1\rangle$ & $\Rmax=0$, costo $0$ y opinión $0$ \\
$E_6$ & 500 & $\langle -100,0.2,3\rangle,\langle 80,0.6,2\rangle,\langle -20,0,1.2\rangle$ & Alcanza para todos \\
$E_7$ & 3   & $\langle 10,0.7,1\rangle,\langle -50,0.5,1\rangle$ & Redondeo de punto flotante \\
$E_8$ & 90  & $\langle 90,0.5,1\rangle,\langle -60,0.5,1\rangle,\langle -60,0.5,1\rangle,\langle 40,0,1\rangle,\langle -70,0.4,2\rangle$ & Óptimos empatados \\
\bottomrule
\end{tabular}
\caption{Entradas diseñadas para analizar el algoritmo voraz.}
\label{tab:entradas}
\end{table}

\begin{table}[h]
\centering
\small
\begin{tabular}{lcccccccc}
\toprule
 & \multicolumn{3}{c}{Voraz} & \multicolumn{3}{c}{Óptimo (FB = PD)} & \\
\cmidrule(lr){2-4}\cmidrule(lr){5-7}
Caso & $E$ & Esf & Ext & $E$ & Esf & Ext & ¿Óptima? \\
\midrule
$E_1$ & \E{1,0,0} & 50 & 33.4996 & \E{1,0,0} & 50 & 33.4996 & Sí \\
$E_2$ & \E{1,0,1} & 28 & 9.4281  & \E{1,0,1} & 28 & 9.4281  & Sí \\
$E_3$ & \E{1,0,0} & 6  & 23.5702 & \E{0,1,1} & 10 & 20.0000 & \textbf{No} \\
$E_4$ & \E{0,1}   & 50 & 10.0000 & \E{0,1}   & 50 & 10.0000 & Sí \\
$E_5$ & \E{1,0,0,0} & 0 & 17.0477 & \E{1,0,0,0} & 0 & 17.0477 & Sí \\
$E_6$ & \E{1,1,1} & 328 & 0.0000 & \E{1,1,1} & 328 & 0.0000 & Sí \\
$E_7$ & \E{1,0}   & 3  & 25.0000 & \E{1,0}   & 3  & 25.0000 & Sí \\
$E_8$ & \E{1,1,0,0,0} & 75 & 24.4949 & \E{1,0,1,0,0} & 75 & 24.4949 & Sí (empate) \\
\bottomrule
\end{tabular}
\caption{Salida del algoritmo voraz frente a la solución óptima.}
\label{tab:voraz}
\end{table}

Comentarios sobre los casos:
\begin{itemize}
  \item $E_3$: costos $6, 5, 5$; beneficios $3600, 2500, 2500$; razones
        $600, 500, 500$. El voraz modera a $a_0$ y le sobran $4$ unidades en
        las que no cabe nadie. El paso 4 empieza también por $a_0$ (mayor
        beneficio). El óptimo modera a $a_1$ y $a_2$ usando todo el
        presupuesto: $B = 5000 > 3600$.
  \item $E_4$: $a_0$ tiene costo $1$ y razón $400$; $a_1$ tiene costo $50$ y
        razón $200$. El paso 3 modera a $a_0$ y ya no cabe $a_1$
        ($B = 400$). El paso 4 empieza por $a_1$ ($B = 10000$) y obtiene el
        óptimo. Sin el paso 4 el extremismo sería $50$ en lugar de $10$.
  \item $E_5$: con $\Rmax = 0$ solo se puede moderar al agente de
        receptividad $1$ (costo $0$). El agente de opinión $0$ no se modera
        porque no aporta beneficio.
  \item $E_7$: sin la tolerancia de \texttt{techo} el costo de $a_0$ sería
        $4 > 3$ y ningún algoritmo lo moderaría.
  \item $E_8$: $\{a_0,a_1\}$ y $\{a_0,a_2\}$ son óptimos (los dos agentes de
        opinión $-60$ son idénticos); los algoritmos devuelven estrategias
        distintas con el mismo esfuerzo y extremismo.
\end{itemize}

\subsection{Complejidad}

Calcular $c_i$ y $b_i$ y filtrar es $O(n)$. Ordenar $|L| \le n$ agentes es
$O(n \log n)$. Cada llamada a \texttt{Llenar} recorre la lista una vez,
$O(n)$, y se hacen a lo sumo dos. Encontrar $b_{\max}$ y calcular el esfuerzo
y el extremismo finales es $O(n)$. En total
\[
  T(n) = O(n) + O(n\log n) + 2\cdot O(n) = \Theta(n \log n),
\]
dominado por el ordenamiento. El espacio es $\Theta(n)$ (la lista ordenada y
dos estrategias).

\subsection{Corrección}
\label{sec:voraz-correccion}

\textbf{No, el algoritmo voraz no siempre da la respuesta correcta.} El caso
$E_3$ es un contraejemplo verificado: el voraz entrega extremismo $23.57$ y el
óptimo es $20$. El motivo es que ModEx es una mochila 0/1, que no tiene la
propiedad de elección voraz: escoger al agente de mejor razón puede dejar
presupuesto sin usar que ningún otro agente alcanza a aprovechar. Ordenar por
razón solo es óptimo en la versión fraccionaria, donde se podría moderar una
fracción de un agente.

\paragraph{Cuándo sí da la respuesta correcta.}
\begin{enumerate}
  \item Si alcanza para todos ($\sum_i c_i \le \Rmax$): modera a todos los
        agentes con $b_i > 0$, lo que da extremismo $0$ (caso $E_6$).
  \item Si $\Rmax = 0$: solo pueden moderarse agentes de costo $0$, y los
        modera a todos (caso $E_5$).
  \item Si todos los costos son iguales: el orden por razón coincide con el
        orden por beneficio y modera a los $\lfloor \Rmax / c\rfloor$ de mayor
        beneficio, que es óptimo.
  \item Si todos los beneficios son iguales: el orden por razón es el de menor
        costo, que maximiza la cantidad de agentes moderados, lo que es
        óptimo.
  \item Si los agentes moderados en el paso 3 son exactamente los primeros de
        la lista ordenada y su costo suma exactamente $\Rmax$: la solución
        alcanza la cota de la mochila fraccionaria y nada puede superarla.
\end{enumerate}

\paragraph{Cuándo no.} Falla cuando, tras moderar a los agentes de mejor
razón, queda presupuesto ``desperdiciado'' que una combinación distinta, de
razón individual menor, sí aprovecharía (caso $E_3$). Esto ocurre con más
frecuencia cuando hay pocos agentes con costos grandes respecto a $\Rmax$.

\begin{proposicion}
Sea $B^*$ el beneficio óptimo y $B_V$ el beneficio de la solución del voraz.
Entonces $B_V \ge B^*/2$.
\end{proposicion}
\begin{proof}
Sean $u_1,\dots,u_m$ los agentes de $L$ en el orden por razón. Si todos caben,
el paso 3 los modera a todos y es óptimo. Si no, sea $u_k$ el primero que no
cabe en el paso 3; los anteriores $u_1,\dots,u_{k-1}$ sí fueron moderados, así
que $B_3 \ge \sum_{l<k} b_{u_l}$. El óptimo de la mochila fraccionaria toma
$u_1,\dots,u_{k-1}$ y una fracción de $u_k$, luego
$B^* \le \sum_{l<k} b_{u_l} + b_{u_k} \le B_3 + b_{\max}$. El paso 4 parte del
agente de beneficio $b_{\max}$, así que $B_4 \ge b_{\max}$ (si ese agente ya
fue moderado en el paso 3, entonces $B_3 \ge b_{\max}$). Por lo tanto
$B_V = \max(B_3, B_4) \ge (B_3 + b_{\max})/2 \ge B^*/2$.
\end{proof}

Sin el paso 4 no hay garantía: con un agente de costo $1$ y beneficio $2$ y
otro de costo $\Rmax$ y beneficio $\Rmax$, el paso 3 solo obtendría $2$
frente a un óptimo de $\Rmax$, arbitrariamente peor (el caso $E_4$ es una
instancia de esta situación). La garantía
es sobre la reducción $B$, no sobre el extremismo: si el óptimo deja la red
casi sin extremismo, el error relativo en $\Ext$ puede ser grande (en las
pruebas aleatorias llegó al $206\%$ en un caso donde el óptimo era cercano a
$0$).

La tabla~\ref{tab:aleatorio} resume el comportamiento observado en
$3600$ redes aleatorias: el voraz entrega el óptimo exacto en la mayoría de
los casos pequeños, y aunque la proporción de óptimos exactos baja al crecer
$n$, el error relativo promedio en el extremismo cae a menos de $0.01\%$
para $n = 1000$, porque con muchos agentes el presupuesto desperdiciado es
pequeño frente al total.

%=====================================================================
\section{Programación dinámica}
%=====================================================================

\subsection{Caracterización de la estructura de una solución óptima}

Consideremos el subproblema $P(i, j)$: \emph{escoger qué agentes moderar entre
los primeros $i$ agentes $a_0,\dots,a_{i-1}$ con un esfuerzo de a lo sumo
$j$, maximizando el beneficio moderado}. El problema original es $P(n,\Rmax)$.

\begin{lema}[Subestructura óptima]
Sea $E^* = \langle e_0,\dots,e_{i-1}\rangle$ una solución óptima de
$P(i,j)$ con $i \ge 1$.
\begin{itemize}
  \item Si $e_{i-1} = 0$, entonces $\langle e_0,\dots,e_{i-2}\rangle$ es una
        solución óptima de $P(i-1, j)$.
  \item Si $e_{i-1} = 1$, entonces $c_{i-1} \le j$ y
        $\langle e_0,\dots,e_{i-2}\rangle$ es una solución óptima de
        $P(i-1, j - c_{i-1})$.
\end{itemize}
\end{lema}
\begin{proof}
Caso $e_{i-1} = 0$: el prefijo usa esfuerzo $\le j$ y solo agentes entre
$a_0,\dots,a_{i-2}$, luego es solución de $P(i-1,j)$. Si existiera una
solución de $P(i-1,j)$ con mayor beneficio, al agregarle $e_{i-1}=0$ se
obtendría una solución de $P(i,j)$ mejor que $E^*$, una contradicción.

Caso $e_{i-1} = 1$: como el esfuerzo total es $\le j$ y $a_{i-1}$ aporta
$c_{i-1}$, el prefijo usa a lo sumo $j - c_{i-1}$. Si existiera una solución
de $P(i-1, j-c_{i-1})$ con mayor beneficio, al agregarle $e_{i-1}=1$ se
tendría una solución de $P(i,j)$ con esfuerzo $\le j$ y beneficio mayor que
el de $E^*$, una contradicción.
\end{proof}

Así, una solución óptima de $P(i,j)$ se construye a partir de una solución
óptima de $P(i-1,j)$ o de $P(i-1,j-c_{i-1})$. Los subproblemas necesarios son
\[
  \{\, P(i,j) : 0 \le i \le n,\; 0 \le j \le A \,\},
  \qquad A = \min\!\left(\Rmax,\ \sum_i c_i\right).
\]
Basta llegar hasta $A$ en lugar de $\Rmax$: con un presupuesto mayor que el
costo de moderar a todos, el presupuesto deja de ser una restricción. Hay
$(n+1)(A+1)$ subproblemas, y como $A \le 300n$, son $O(n^2)$.

\subsection{Definición recursiva del valor de una solución óptima}

Sea $M(i,j)$ el máximo beneficio de $P(i,j)$:
\[
  M(i,j) =
  \begin{cases}
    0 & \text{si } i = 0,\\[2pt]
    M(i-1, j) & \text{si } i > 0 \text{ y } c_{i-1} > j,\\[2pt]
    \max\bigl( M(i-1,j),\ M(i-1, j - c_{i-1}) + b_{i-1} \bigr) & \text{si } i > 0 \text{ y } c_{i-1} \le j.
  \end{cases}
\]
El costo (extremismo) de una solución óptima del problema original es
\[
  \Ext^* = \frac{\sqrt{S - M(n, A)}}{n}, \qquad S = \sum_i b_i .
\]
De manera equivalente, el extremismo restante mínimo
$X(i,j) = \sum_{l<i} b_l - M(i,j)$ satisface $X(0,j) = 0$,
$X(i,j) = X(i-1,j) + b_{i-1}$ si $c_{i-1} > j$ y
$X(i,j) = \min\bigl(X(i-1,j) + b_{i-1},\ X(i-1,j-c_{i-1})\bigr)$ en otro caso,
con $\Ext^* = \sqrt{X(n,A)}/n$.

\subsection{Algoritmo para calcular el costo de una solución óptima}

La fila $i$ de $M$ solo depende de la fila $i-1$, así que la tabla se llena
de abajo hacia arriba, fila por fila:

\begin{algorithm}[h]
\caption{Costo óptimo por programación dinámica}
\label{alg:pd-costo}
\KwIn{Red $RS$ con $n$ agentes}
\KwOut{$\Ext^*$}
calcular $c_i$ y $b_i$ para $0 \le i < n$; $S \gets \sum_i b_i$\;
$A \gets \min(\Rmax, \sum_i c_i)$\;
\lFor{$j \gets 0$ \KwTo $A$}{$M[0][j] \gets 0$}
\For{$i \gets 1$ \KwTo $n$}{
  \For{$j \gets 0$ \KwTo $A$}{
    $M[i][j] \gets M[i-1][j]$\;
    \If{$c_{i-1} \le j$ \textbf{y} $M[i-1][j-c_{i-1}] + b_{i-1} > M[i][j]$}{
      $M[i][j] \gets M[i-1][j-c_{i-1}] + b_{i-1}$\;
    }
  }
}
\Return $\sqrt{S - M[n][A]}\,/\,n$\;
\end{algorithm}

\subsection{Algoritmo para calcular una solución óptima}

Se agrega una tabla de decisiones $D[i][j] \in \{0,1\}$ con $D[i][j] = 1$ si
y solo si el valor $M[i][j]$ se obtuvo moderando a $a_{i-1}$ (es decir, si en
el algoritmo~\ref{alg:pd-costo} se ejecutó la línea de la condición). Luego
se recorre la tabla desde $(n, A)$ hacia atrás:

\begin{algorithm}[h]
\caption{Reconstrucción de una estrategia óptima}
\KwIn{Tabla $D$, costos $c_i$, $n$, $A$}
\KwOut{Estrategia óptima $E$}
$E \gets \langle 0,\dots,0\rangle$; $j \gets A$\;
\For{$i \gets n$ \textbf{bajando hasta} $1$}{
  \If{$D[i][j] = 1$}{
    $e_{i-1} \gets 1$; $j \gets j - c_{i-1}$\;
  }
}
\Return $(E, \Esf(RS,E), \Ext(\Mod(RS,E)))$\;
\end{algorithm}

En cada paso, si $D[i][j] = 1$ el óptimo de $P(i,j)$ modera a $a_{i-1}$ y
continúa en $P(i-1, j - c_{i-1})$; si no, continúa en $P(i-1, j)$. Por el
lema de subestructura óptima, la estrategia construida tiene beneficio
$M(n,A)$ y es óptima. La reconstrucción toma $O(n)$.

\paragraph{En la implementación.} Como la fila $i$ solo depende de la fila
$i-1$, \texttt{modexPD} guarda los valores en un solo arreglo de $A+1$ reales
que se recorre con $j$ de $A$ hacia $0$ (así $M[j - c]$ todavía corresponde a
la fila anterior), y la tabla $D$ como una matriz de bits
(\texttt{std::vector<bool>}). El núcleo del algoritmo es:

\begin{lstlisting}
std::vector<double> M(A + 1, 0.0);
std::vector<std::vector<bool>> moderar(n, std::vector<bool>(A + 1, false));
for (std::size_t i = 0; i < n; ++i) {
    const auto c = static_cast<std::size_t>(costo[i]);
    if (c > A) continue;
    for (std::size_t j = A + 1; j-- > c;) {
        double conAgente = M[j - c] + beneficio[i];
        if (conAgente > M[j]) {
            M[j] = conAgente;
            moderar[i][j] = true;
        }
    }
}
std::vector<int> E(n, 0);
std::size_t j = A;
for (std::size_t i = n; i-- > 0;) {
    if (moderar[i][j]) {
        E[i] = 1;
        j -= static_cast<std::size_t>(costo[i]);
    }
}
\end{lstlisting}

\subsection{Complejidad}

\paragraph{Complejidad en tiempo.} Calcular costos y beneficios es $O(n)$.
Llenar la tabla son $(n+1)(A+1)$ celdas de costo constante, y la
reconstrucción es $O(n)$. Por lo tanto $T(n) = \Theta(n \cdot A)$. Como
$A \le \sum_i c_i \le 300n$,
\[
  T(n) = O(300\, n^2) = O(n^2).
\]
En general el algoritmo es pseudopolinomial en $\Rmax$, pero en este problema
el presupuesto útil está acotado por $300n$, lo que lo hace polinomial en $n$.

\paragraph{Complejidad en espacio.} Las tablas $M$ y $D$ tienen
$(n+1)(A+1) \approx n\cdot A \le 300n^2$ celdas, luego el espacio es
$\Theta(n\cdot A) = O(n^2)$. (La implementación reduce $M$ a una fila, pero
$D$ sigue siendo de $n \times (A+1)$, por lo que el orden se mantiene.)

\paragraph{¿Es útil en la práctica?} Supongamos un equipo que procesa
$3\times 10^8$ operaciones por minuto y usa $4$ bytes por celda, con
$n = 2^k$. Tomamos el peor caso $A = 300n$ (con las aproximaciones
$n+1\sim n$, $A+1 \sim A$):
\[
  \text{operaciones} \approx n\cdot A = 300 n^2 = 300\cdot 4^k
  \;\Rightarrow\;
  \text{tiempo} \approx \frac{300\cdot 4^k}{3\times 10^8} = \frac{4^k}{10^6}\ \text{minutos},
\]
\[
  \text{espacio} \approx 4 \cdot 300 n^2 = 1200 \cdot 2^{2k}\ \text{bytes}.
\]
Con $1\ \text{año} = 518400$ minutos y $2^{10}$ bytes $= 1$ KB, etc.:

\begin{table}[h]
\centering
\begin{tabular}{rrrrr}
\toprule
$k$ & $n = 2^k$ & Tiempo & Espacio & Veredicto \\
\midrule
5  & 32        & $0.001$ min ($0.06$ s)        & $1.17$ MB  & Aceptable \\
8  & 256       & $0.066$ min ($3.9$ s)         & $75$ MB    & Aceptable \\
10 & 1\,024    & $1.05$ min                    & $1.17$ GB  & Aceptable \\
12 & 4\,096    & $16.8$ min                    & $18.75$ GB & Tiempo aceptable, espacio no \\
14 & 16\,384   & $268$ min ($4.5$ h)           & $300$ GB   & Inaceptable \\
16 & 65\,536   & $4\,295$ min ($3$ días)       & $4.69$ TB  & Inaceptable \\
18 & 262\,144  & $68\,719$ min ($0.13$ años)   & $75$ TB    & Inaceptable \\
20 & 1\,048\,576 & $1.1\times 10^6$ min ($2.1$ años) & $1.17$ PB & Inaceptable \\
\bottomrule
\end{tabular}
\caption{Tiempo y espacio de la programación dinámica en el peor caso ($A = 300n$).}
\label{tab:pd-practica}
\end{table}

\begin{itemize}
  \item \textbf{Tiempo:} es aceptable hasta $k \approx 12$ (unos 17 minutos);
        $k = 14$ ya toma más de 4 horas y para $k \ge 16$ toma días o años.
  \item \textbf{Espacio:} es aceptable hasta $k \approx 10$ (un poco más de
        1 GB, que cabe en la memoria de un computador personal); para
        $k = 12$ se necesitan casi 19 GB y para $k \ge 14$ cientos de GB o
        más, lo que es inaceptable. El espacio es el factor que limita
        primero.
\end{itemize}

Este es el peor caso, con todos los agentes extremos, poco receptivos y muy
influyentes, y un $\Rmax$ que alcanza para moderarlos a todos. En las redes
aleatorias de las pruebas, con $\Rmax$ igual a la mitad del esfuerzo total,
se obtuvo $A \approx 25n$, doce veces menos. Además, la implementación guarda
$D$ con 1 bit por celda en lugar de 4 bytes, lo que divide la memoria entre
32: para $k = 12$ en el peor caso usaría unos $600$ MB. El equipo de prueba
es también mucho más rápido que el modelo del enunciado: llena unas
$8\times 10^8$ celdas por segundo (tabla~\ref{tab:tiempos}).

%=====================================================================
\section{Implementación}
%=====================================================================

\subsection{Tecnología y estructura}

El programa está escrito en C++17 con Qt 6 (Widgets) para la interfaz y
CMake para la construcción. La lógica está en una librería
(\texttt{modex\_core}) que usan tanto la aplicación gráfica (\texttt{ModEx})
como un programa de pruebas en consola (\texttt{pruebas}).

\begin{table}[h]
\centering
\small
\begin{tabular}{ll}
\toprule
Archivo & Contenido \\
\midrule
\texttt{include/agente.hpp}, \texttt{src/agente.cpp} & Clase \texttt{Agente} y generación aleatoria \\
\texttt{include/red.hpp}, \texttt{src/red.cpp} & Clase \texttt{Red}, $\Ext$, costo, esfuerzo, archivos \\
\texttt{include/resultado.hpp} & Tripleta $(E, \Esf, \Ext)$ \\
\texttt{src/estrategias/fuerzabruta.cpp} & \texttt{modexFB} \\
\texttt{src/estrategias/voraz.cpp} & \texttt{modexV} \\
\texttt{src/estrategias/dinamica.cpp} & \texttt{modexPD} \\
\texttt{src/mainwindow.cpp} & Interfaz gráfica \\
\texttt{pruebas/pruebas.cpp} & Pruebas automáticas y mediciones \\
\texttt{casos/} & Ejemplos del enunciado en el formato de entrada \\
\bottomrule
\end{tabular}
\caption{Archivos principales del proyecto.}
\end{table}

Las tres funciones pedidas, \texttt{modexFB}, \texttt{modexV} y
\texttt{modexPD}, reciben una red (\texttt{const Red\&}) y devuelven un
\texttt{Resultado} con la estrategia $E$, su esfuerzo y el extremismo de la
red moderada.

\subsection{Interfaz}

La interfaz permite:
\begin{itemize}
  \item \textbf{Abrir} un archivo de entrada y visualizar la red: una tabla
        con la opinión, la receptividad, la influencia y el costo de moderar
        a cada agente, además de $n$, $\Rmax$, el esfuerzo de moderar a todos
        y el extremismo inicial.
  \item \textbf{Generar casos de prueba} aleatorios con $n$ y $\Rmax$
        escogidos por el usuario; el caso se guarda en un archivo con el
        formato de entrada y se carga.
  \item Ejecutar \textbf{Fuerza bruta}, \textbf{Voraz} o \textbf{Dinámica}, o
        \textbf{Comparar} las tres (la fuerza bruta solo se habilita para
        $n \le 25$). Los algoritmos corren en un hilo aparte para no congelar
        la ventana. La estrategia de cada algoritmo aparece en su propia
        columna de la tabla de agentes.
  \item \textbf{Revisar la calidad de las salidas}: una tabla de resultados
        muestra para cada algoritmo $\Ext$, $\Esf$, si $\Esf \le \Rmax$, si
        el recálculo de $\Esf$ y $\Ext$ a partir de $E$ coincide con lo
        reportado, el tiempo de ejecución y si la solución es óptima
        comparada con la de programación dinámica (y por cuánto difiere si
        no lo es). Con la columna de costos es fácil comprobar a mano el
        esfuerzo de una solución.
  \item \textbf{Guardar la salida} del algoritmo seleccionado en un archivo
        de texto con el formato pedido.
\end{itemize}

\subsection{Formatos de entrada y salida}

La entrada tiene $n+2$ líneas: $n$; una línea
\texttt{op,recep,inf} por agente; y $\Rmax$. El lector acepta como
separadores comas, punto y coma, espacios y ``\texttt{.~}'' (punto seguido de
espacio, como aparece en el enunciado), tolera líneas vacías y finales de
línea de Windows, y valida los rangos ($o_i \in [-100,100]$,
$r_i \in [0,1]$, $inf_i \in [1,3]$, $\Rmax \ge 0$), informando la línea del
error. La salida tiene $n+2$ líneas: $\Ext$ (con seis decimales), $\Esf$, y
un $0$ o $1$ por agente.

%=====================================================================
\section{Comparación de las tres soluciones}
%=====================================================================

\subsection{Diseño de los casos de prueba}

Los casos se diseñaron en cuatro grupos:
\begin{enumerate}
  \item \textbf{Ejemplos del enunciado} ($E_1$, $E_2$), para validar el
        cálculo del esfuerzo y del extremismo contra valores conocidos.
  \item \textbf{Casos borde} ($E_5$, $E_6$, $E_7$): $\Rmax = 0$, presupuesto
        que alcanza para todos, agentes con receptividad $1$ (costo $0$) y
        opinión $0$ (beneficio $0$), y valores que exponen errores de
        redondeo.
  \item \textbf{Casos adversarios para el voraz} ($E_3$, $E_4$, $E_8$),
        construidos a partir del análisis de la sección~\ref{sec:voraz-correccion}:
        presupuesto desperdiciado, agente muy extremo bloqueado por uno
        barato y óptimos empatados.
  \item \textbf{Casos aleatorios}, generados con la misma función que usa la
        interfaz: $o_i$ uniforme en $\{-100,\dots,100\}$, $r_i$ uniforme en
        $[0,1]$ e $inf_i$ uniforme en $[1,3]$ (con dos decimales), y $\Rmax$
        uniforme entre $0$ y $\sum_i c_i$ para cubrir desde presupuestos muy
        escasos hasta suficientes. Se usó una semilla fija para que los
        resultados sean reproducibles.
\end{enumerate}

Para cada caso se verifica automáticamente (programa \texttt{pruebas}) que
las tres salidas sean aplicables, que el esfuerzo y el extremismo reportados
coincidan con los recalculados a partir de $E$, y que fuerza bruta y
programación dinámica den el mismo extremismo. Todas las verificaciones
pasaron; en particular, \textbf{la programación dinámica coincidió con la
fuerza bruta en los $2000$ casos aleatorios con $n \le 14$} del programa de
pruebas y en los $3200$ casos con $n \le 20$ de la tabla~\ref{tab:aleatorio}.

\subsection{Resultados en los casos diseñados}

\begin{table}[h]
\centering
\small
\begin{tabular}{lrrrrrrr}
\toprule
 & & \multicolumn{2}{c}{Fuerza bruta} & \multicolumn{2}{c}{Voraz} & \multicolumn{2}{c}{Dinámica} \\
\cmidrule(lr){3-4}\cmidrule(lr){5-6}\cmidrule(lr){7-8}
Caso & Ext inicial & Esf & Ext & Esf & Ext & Esf & Ext \\
\midrule
$E_1$ & 47.2582 & 50  & 33.4996 & 50  & 33.4996 & 50  & 33.4996 \\
$E_2$ & 21.6025 & 28  & 9.4281  & 28  & 9.4281  & 28  & 9.4281 \\
$E_3$ & 30.9121 & 10  & 20.0000 & 6   & \textbf{23.5702} & 10 & 20.0000 \\
$E_4$ & 50.9902 & 50  & 10.0000 & 50  & 10.0000 & 50  & 10.0000 \\
$E_5$ & 33.0246 & 0   & 17.0477 & 0   & 17.0477 & 0   & 17.0477 \\
$E_6$ & 69.3462 & 328 & 0.0000  & 328 & 0.0000  & 328 & 0.0000 \\
$E_7$ & 25.4951 & 3   & 25.0000 & 3   & 25.0000 & 3   & 25.0000 \\
$E_8$ & 32.6803 & 75  & 24.4949 & 75  & 24.4949 & 75  & 24.4949 \\
\bottomrule
\end{tabular}
\caption{Resultados de los tres algoritmos en los casos diseñados.}
\label{tab:comparacion}
\end{table}

\subsection{Calidad del voraz en casos aleatorios}

Para $n \le 20$ el óptimo se calculó con fuerza bruta (y se comprobó que la
dinámica da lo mismo); para $n = 100$ y $n = 1000$ se usó la dinámica. El
error relativo es $(\Ext_V - \Ext^*)/\Ext^*$, y $B_V/B^*$ es la fracción del
beneficio óptimo que logra el voraz.

\begin{table}[h]
\centering
\begin{tabular}{rrrrrr}
\toprule
$n$ & Casos & Voraz óptimo & Error prom. en Ext & Error máx. en Ext & $B_V/B^*$ mín. \\
\midrule
5    & 1000 & 88.4\% & 2.47\%  & 206.2\% & 0.729 \\
10   & 1000 & 71.9\% & 2.75\%  & 105.8\% & 0.794 \\
15   & 1000 & 61.6\% & 2.24\%  & 147.2\% & 0.778 \\
20   & 200  & 56.0\% & 1.97\%  & 60.2\%  & 0.888 \\
100  & 200  & 27.5\% & 0.21\%  & 3.09\%  & 0.982 \\
1000 & 200  & 17.5\% & 0.005\% & 0.12\%  & 0.997 \\
\bottomrule
\end{tabular}
\caption{Calidad del algoritmo voraz en redes aleatorias.}
\label{tab:aleatorio}
\end{table}

Se observa que:
\begin{itemize}
  \item El voraz siempre logró al menos el $72.9\%$ del beneficio óptimo,
        por encima de la garantía teórica del $50\%$.
  \item La probabilidad de obtener el óptimo exacto baja al crecer $n$, pero
        el error relativo cae rápidamente: para $n = 1000$ el extremismo del
        voraz difiere en promedio un $0.005\%$ del óptimo.
  \item Los errores máximos grandes en $\Ext$ aparecen en redes pequeñas
        donde el óptimo deja un extremismo cercano a $0$, de modo que una
        diferencia absoluta pequeña es grande en términos relativos.
\end{itemize}

\subsection{Tiempos de ejecución}

Las mediciones se hicieron en un Intel Core i7-1165G7 a 2.8 GHz con
compilación optimizada (\texttt{-O2}), con $\Rmax$ igual a la mitad del
esfuerzo de moderar a todos.

\begin{table}[h]
\centering
\begin{minipage}{0.38\textwidth}
\centering
\begin{tabular}{rr}
\toprule
$n$ & FB (ms) \\
\midrule
4  & 0.002 \\
8  & 0.035 \\
12 & 0.746 \\
16 & 12.9 \\
18 & 61.0 \\
20 & 296.7 \\
22 & 1\,203.3 \\
24 & 4\,403.7 \\
\bottomrule
\end{tabular}
\end{minipage}%
\begin{minipage}{0.58\textwidth}
\centering
\begin{tabular}{rrrr}
\toprule
$n$ & $A$ & Voraz (ms) & Dinámica (ms) \\
\midrule
10   & 267       & 0.0006 & 0.005 \\
20   & 554       & 0.0007 & 0.020 \\
100  & 2\,604    & 0.0034 & 0.417 \\
500  & 11\,874   & 0.018  & 8.61 \\
1000 & 26\,583   & 0.059  & 36.7 \\
2000 & 52\,536   & 0.142  & 130.2 \\
5000 & 126\,149  & 0.450  & 788.6 \\
\bottomrule
\end{tabular}
\end{minipage}
\caption{Tiempos de ejecución medidos.}
\label{tab:tiempos}
\end{table}

\begin{figure}[h]
\centering
\begin{tikzpicture}
\begin{loglogaxis}[
  width=0.85\textwidth, height=7cm,
  xlabel={$n$}, ylabel={tiempo (ms)},
  legend pos=north west, grid=major,
  xmin=3, xmax=8000
]
\addplot[mark=*,red] coordinates {
  (4,0.0022) (8,0.0347) (12,0.746) (16,12.86) (18,60.97) (20,296.7) (22,1203.3) (24,4403.7)
};
\addlegendentry{Fuerza bruta}
\addplot[mark=square*,blue] coordinates {
  (10,0.0049) (20,0.0197) (100,0.4169) (500,8.6118) (1000,36.652) (2000,130.17) (5000,788.61)
};
\addlegendentry{Dinámica}
\addplot[mark=triangle*,green!50!black] coordinates {
  (10,0.0006) (20,0.0007) (100,0.0034) (500,0.0181) (1000,0.059) (2000,0.1424) (5000,0.4503)
};
\addlegendentry{Voraz}
\end{loglogaxis}
\end{tikzpicture}
\caption{Tiempos de ejecución en escala logarítmica.}
\label{fig:tiempos}
\end{figure}

Los tiempos confirman el análisis: la fuerza bruta se multiplica por algo
más de $2$ con cada agente adicional ($\Theta(n\cdot 2^n)$); la dinámica
crece aproximadamente con el cuadrado de $n$ (al pasar de $n=1000$ a
$n=2000$ el tiempo se multiplica por $3.6$, y $A$ crece linealmente con $n$);
el voraz es casi lineal. Extrapolando, la fuerza bruta tardaría unos $6$
minutos con $n = 30$, unos $6$ días con $n = 40$ y unos $20$ años con $n = 50$,
mientras que la dinámica resuelve $n = 5000$ en menos de un segundo.

%=====================================================================
\section{Conclusiones}
%=====================================================================

\begin{table}[h]
\centering
\begin{tabular}{lccc}
\toprule
 & Fuerza bruta & Voraz & Dinámica \\
\midrule
Tiempo  & $\Theta(n\cdot 2^n)$ & $\Theta(n\log n)$ & $\Theta(n\cdot A) = O(n^2)$ \\
Espacio & $\Theta(n)$ & $\Theta(n)$ & $\Theta(n\cdot A) = O(n^2)$ \\
¿Óptimo? & Siempre & No siempre ($\ge 1/2$ de $B^*$) & Siempre \\
Tamaño práctico & $n \lesssim 25$ & Millones de agentes & Miles de agentes \\
\bottomrule
\end{tabular}
\caption{Resumen de las tres estrategias.}
\end{table}

\begin{itemize}
  \item \textbf{Fuerza bruta.} Su ventaja es la simplicidad y la certeza: es
        fácil de implementar y de demostrar correcta, y sirve como referencia
        para validar los otros algoritmos (así se usó en las pruebas). Su
        desventaja es el crecimiento exponencial, que la hace inútil más allá
        de unos 25 agentes.
  \item \textbf{Voraz.} Es, por mucho, el más rápido, con un costo dominado
        por un ordenamiento, y usa memoria lineal. En la práctica sus
        soluciones son muy buenas, y con muchos agentes casi indistinguibles
        del óptimo. Su desventaja es que no garantiza el óptimo: existen
        entradas pequeñas donde se equivoca, y solo garantiza la mitad del
        beneficio óptimo. Es la mejor opción cuando la red es muy grande o
        se necesita una respuesta inmediata y se tolera un error pequeño.
  \item \textbf{Programación dinámica.} Garantiza el óptimo en tiempo
        polinomial gracias a que el presupuesto útil está acotado por $300n$,
        y resolvió miles de agentes en menos de un segundo. Su desventaja es
        la memoria: la tabla de decisiones crece como $n^2$ y es el primer
        recurso que se agota (alrededor de $n \approx 2^{10}$ a $2^{12}$ con
        4 bytes por celda). Es la mejor opción cuando se necesita la solución
        exacta y la red tiene hasta algunos miles de agentes.
  \item Reconocer que ModEx es una mochila 0/1 fue clave: explica por qué el
        voraz por razón falla, de dónde sale su garantía, y cuál es la
        subestructura óptima que aprovecha la programación dinámica.
\end{itemize}

\end{document}
